\documentclass[11pt]{article}
\usepackage[margin=1in]{geometry}
\usepackage{amsmath,amssymb,amsthm,mathtools}
\usepackage{graphicx}
\usepackage{microtype}
\usepackage[hidelinks]{hyperref}

\newtheorem{theorem}{Theorem}
\newtheorem{lemma}[theorem]{Lemma}
\newcommand{\PP}{\mathsf P}
\newcommand{\DD}{\mathsf D}
\newcommand{\TT}{\mathsf T}
\newcommand{\CC}{\mathcal C}

\title{All-aperture self-adjointness for the MIT bag model\\
on an unbounded circular cone}
\author{Candidate full solution packet for arXiv:2201.08192}
\date{12 August 2026}

\begin{document}
\maketitle

\begin{abstract}
Cassano and Lotoreichik proved essential self-adjointness of the MIT bag Dirac
operator on convex circular cones and conjectured it for every aperture.  Their
reduction leaves one Ferrers-function equation for the re-entrant range.  We
prove that the left side minus the right side of this equation is strictly
positive throughout the critical spectral interval.  The Mehler--Dirichlet
formula turns the expression into an integral which, for nonnegative spectral
parameter, is minimized at the endpoint $\lambda=1/2$.  That endpoint reduces
to a strict elementary inequality between complete elliptic integrals.  This
closes the last spectral case and proves the conjecture for all
$0<\omega<\pi$.
\end{abstract}

\section{Source conjecture}

The source is B. Cassano and V. Lotoreichik, ``Self-adjointness for the MIT bag
model on an unbounded cone,'' arXiv:2201.08192, later published in
\emph{Mathematische Nachrichten} 297 (2024), 1006--1041
\cite{CassanoLotoreichik}.  On physical PDF page~6, Remark~2.4 states:

\begin{quote}
``We conjecture that Theorem 2.3 is true for all the values
$\omega\in(0,\pi)$ of the aperture angle. In Appendix C we provide numerical
evidence to support our conjecture, showing that for
$\pi/2<\omega<\pi$ all the eigenvalues of the spin-orbit operator $K_\omega$
escape the critical interval $[-1/2,1/2]$. This implies essential
self-adjointness of $D_\omega$ and the validity of (2.9) for all
$\omega\in(0,\pi)$.''
\end{quote}

The rendered source-page evidence, including the reduction immediately before
the conjecture, is reproduced below.

\begin{center}
\includegraphics[width=0.97\textwidth]{figures/open_problem_crop.png}
\end{center}

Here $\DD_\omega$ is the massless Dirac operator on the circular cone
$\CC_\omega$ with MIT bag boundary conditions, initially defined away from
the vertex.  The conjectured domain of its closure is
\begin{equation}\label{eq:domain}
 \operatorname{dom}\overline{\DD_\omega}
 =\left\{u\in H^1(\CC_\omega;\mathbb C^4):
 u|_{\partial\CC_\omega}
 +i\beta(\alpha\!\cdot\!\nu_\omega)
 u|_{\partial\CC_\omega}=0\right\}.
\end{equation}

\section{The exact missing scalar statement}

Write $\PP_\nu^\mu(x)$ for the Ferrers function of the first kind and
$\PP_\nu=\PP_\nu^0$.  The source proves the conjecture for
$0<\omega<\pi/2$ and records that the half-space case $\omega=\pi/2$ is
known.  For $\pi/2<\omega<\pi$, Appendix~C makes the remaining reduction
explicit:
\begin{itemize}
\item all angular modes $k\notin\{-1,0\}$ already have
      $\sigma(\TT_k)\cap[-1/2,1/2]=\varnothing$;
\item the $k=-1$ and $k=0$ cases are equivalent;
\item the relevant $k=0$ set consists of the roots of
\begin{equation}\label{eq:source-equation}
 (\lambda+1)\PP_\lambda^{-1}(\cos\omega)
 =\PP_{\lambda-1}(\cos\omega).
\end{equation}
\end{itemize}
Thus it suffices to prove that \eqref{eq:source-equation} has no root in
$[-1/2,1/2]$ for a re-entrant aperture.

\section{Proof intuition}

The source equation looks delicate because the two Ferrers functions can each
change sign near $x=-1$.  A recurrence combines them before any sign estimate
is made.  The Mehler--Dirichlet formula then writes the combined expression as
an integral whose dependence on $\lambda$ is elementary trigonometry.
Negative $\lambda$ gives a pointwise positive integrand.  For
$0\le\lambda\le1/2$, that integrand decreases with $\lambda$, so the entire
interval reduces to $\lambda=1/2$.  A half-angle substitution evaluates the
endpoint using complete elliptic integrals $E$ and $K$.  Pairing the integrands
at $y$ and $\pi/2-y$ proves the needed strict inequality $E>qK$.

\section{Ferrers positivity lemma}

\begin{lemma}\label{lem:positive}
For every $\theta\in(\pi/2,\pi)$ and every
$\lambda\in[-1/2,1/2]$,
\begin{equation}\label{eq:F}
 F_\theta(\lambda):=
 (\lambda+1)\PP_\lambda^{-1}(\cos\theta)
 -\PP_{\lambda-1}(\cos\theta)>0.
\end{equation}
\end{lemma}

\begin{proof}
Put
\[
 x=\cos\theta<0,\qquad s=\sin\theta>0,
 \qquad p=\sin\frac\theta2,\qquad q=\cos\frac\theta2.
\]
Then $p>q>0$.  The degree symmetry $\PP_{\lambda-1}=\PP_{-\lambda}$ and a
standard Ferrers recurrence give
\begin{equation}\label{eq:recurrence}
 (\lambda+1)\PP_\lambda^{-1}(x)
 =\frac{\PP_{-\lambda}(x)-x\PP_\lambda(x)}{s}.
\end{equation}
The identity is regular also at $\lambda=0$.  Consequently,
\begin{equation}\label{eq:H}
 sF_\theta(\lambda)
 =(1-s)\PP_{-\lambda}(x)-x\PP_\lambda(x)
 =:H_\theta(\lambda).
\end{equation}

For $0<\theta<\pi$, the Mehler--Dirichlet formula is
\begin{equation}\label{eq:MD}
 \PP_\nu(\cos\theta)
 =\frac{\sqrt2}{\pi}\int_0^\theta
 \frac{\cos((\nu+\tfrac12)t)}
 {\sqrt{\cos t-\cos\theta}}\,dt
 \qquad\text{\cite[Eq.~14.12.1]{DLMF}}.
\end{equation}
Using
\[
 1-s-x=2p(p-q),\qquad 1-s+x=-2q(p-q),
\]
and expanding the two cosines in \eqref{eq:MD}, equation \eqref{eq:H}
becomes
\begin{equation}\label{eq:Hintegral}
 H_\theta(\lambda)=\frac{2\sqrt2(p-q)}{\pi}
 \int_0^\theta
 \frac{J_\lambda(t)}{\sqrt{\cos t-\cos\theta}}\,dt,
\end{equation}
where
\begin{equation}\label{eq:J}
 J_\lambda(t)=
 p\cos\frac t2\cos(\lambda t)
 -q\sin\frac t2\sin(\lambda t).
\end{equation}

If $-1/2\le\lambda\le0$, then
$\cos(\lambda t)>0$ and $\sin(\lambda t)\le0$ for $0<t<\theta<\pi$.
Hence $J_\lambda(t)>0$ pointwise, and \eqref{eq:Hintegral} is positive.

Suppose now that $0\le\lambda\le1/2$.  For $0<t<\theta$,
\begin{equation}\label{eq:monotone}
 \frac{\partial J_\lambda(t)}{\partial\lambda}
 =-t\left[
 p\cos\frac t2\sin(\lambda t)
 +q\sin\frac t2\cos(\lambda t)
 \right]\le0.
\end{equation}
Therefore $J_\lambda(t)\ge J_{1/2}(t)$ and it remains only to prove that the
integral of
\[
 J_{1/2}(t)=p\cos^2\frac t2-q\sin^2\frac t2
\]
is positive.

Let $u=\sin(t/2)/p$.  With the standard complete elliptic integrals
\[
 K(p)=\int_0^{\pi/2}\frac{dy}{\sqrt{1-p^2\sin^2y}},
 \qquad
 E(p)=\int_0^{\pi/2}\sqrt{1-p^2\sin^2y}\,dy,
\]
the elementary substitution, followed by $u=\sin y$, gives
\begin{align}
 \int_0^\theta\frac{J_{1/2}(t)}
 {\sqrt{\cos t-\cos\theta}}\,dt
 &=\sqrt2\int_0^1
 \frac{p-p^2(p+q)u^2}
 {\sqrt{1-u^2}\sqrt{1-p^2u^2}}\,du \notag\\
 &=\sqrt2\big[(p+q)E(p)-qK(p)\big]. \label{eq:elliptic}
\end{align}

We finish with a short proof that $E(p)>qK(p)$.  Set
\[
 a(y)=\sqrt{1-p^2\sin^2y},\qquad
 b(y)=a(\tfrac\pi2-y)=\sqrt{1-p^2\cos^2y}.
\]
Because $p^2+q^2=1$,
\[
 a(y)^2b(y)^2=q^2+p^4\sin^2y\cos^2y\ge q^2.
\]
Pairing $y$ with $\pi/2-y$ therefore yields
\begin{equation}\label{eq:paired}
 \left(a-\frac q a\right)+\left(b-\frac q b\right)
 =\frac{(a+b)(ab-q)}{ab}\ge0,
\end{equation}
strictly for almost every $y$.  Integrating proves $E(p)-qK(p)>0$.
Since $p+q>1$, the bracket in \eqref{eq:elliptic} is strictly positive.
Equations \eqref{eq:Hintegral}--\eqref{eq:elliptic} show
$H_\theta(1/2)>0$, hence $H_\theta(\lambda)>0$ for
$0\le\lambda\le1/2$.  Together with the negative-$\lambda$ case and
$s>0$, this proves \eqref{eq:F}.
\end{proof}

\section{Full resolution of the conjecture}

\begin{theorem}\label{thm:main}
For every half-aperture $\omega\in(0,\pi)$, the MIT bag Dirac operator
$\DD_\omega$ of Cassano--Lotoreichik is essentially self-adjoint.  Its closure
has exactly the $H^1$ domain in \eqref{eq:domain}.
\end{theorem}

\begin{proof}
For $0<\omega<\pi/2$, this is Theorem~2.3 of the source
\cite{CassanoLotoreichik}; for $\omega=\pi/2$, the half-space result recalled
immediately after its Theorem~2.1 applies.

Let $\pi/2<\omega<\pi$.  The source proves that
$\sigma(\TT_k)\cap[-1/2,1/2]=\varnothing$ for every
$k\notin\{-1,0\}$.  It also proves that the two remaining modes are equivalent
and that the $k=0$ roots are precisely the solutions of
\eqref{eq:source-equation}.  Lemma~\ref{lem:positive} shows that no such
solution lies in $[-1/2,1/2]$.  Hence
\[
 \sigma(\TT_k)\cap[-1/2,1/2]=\varnothing
 \qquad(k\in\mathbb Z).
\]
The orthogonal fiber decomposition and one-dimensional limit-point criterion
in Theorems~2.1 and 2.3 of the source now give essential self-adjointness and
the domain formula \eqref{eq:domain}.  This is exactly the implication stated
in Remark~2.4 of the source.
\end{proof}

\section{Verification, novelty, and limitations}

The proof is analytic and does not rely on computation.  The included
high-precision verifier independently checks the Ferrers recurrence reduction,
the Mehler--Dirichlet integral, the monotonic endpoint reduction, and the
elliptic endpoint identity.  It evaluates the source expression on 1,312
$(\omega,\lambda)$ pairs, including apertures within $10^{-6}$ of $\pi/2$ and
near $\pi$, and returns PASS.  See \texttt{VERIFIER\_REPORT.md}.

A bounded novelty search on 12 August 2026 covered the four run indexes,
arXiv, the published article page, exact-title and exact-conjecture searches,
the author pair, and combinations of ``MIT bag,'' ``non-convex cone,''
``spin-orbit operator,'' and ``critical interval.''  It found the source and
its published version, but no later paper resolving the all-aperture
conjecture or proving the Ferrers inequality above.  This is not a guarantee
of priority.

The result resolves the circular-cone conjecture exactly as stated.  It does
not address general non-circular cones mentioned in Remark~2.5, nor does it
extend the perturbative Hardy estimate of the source from convex to
re-entrant cones.  Human review should focus on the source-to-scalar reduction
and the coefficient algebra leading from \eqref{eq:MD} to
\eqref{eq:Hintegral}; the remaining endpoint argument is elementary.

\begin{thebibliography}{9}
\bibitem{CassanoLotoreichik}
B. Cassano and V. Lotoreichik,
\emph{Self-adjointness for the MIT bag model on an unbounded cone},
Math. Nachr. \textbf{297} (2024), 1006--1041; arXiv:2201.08192.

\bibitem{DLMF}
T. M. Dunster, \emph{Legendre and Related Functions}, Chapter~14 in the
NIST Digital Library of Mathematical Functions, especially
Sections~14.9, 14.10, and 14.12,
\url{https://dlmf.nist.gov/14}.
\end{thebibliography}

\end{document}
